\chapter{Contexte et objectifs}
\label{chap:contexte}

\section{Problématique clinique}

Le syndrome de défilé cervico-thoracique (\emph{Thoracic Outlet Syndrome}, TOS) est une pathologie
liée à la compression de structures neurovasculaires au niveau du défilé cervico-thoracique lors de
certains mouvements du bras. Le diagnostic s'appuie sur une IRM \textbf{dynamique} : une séquence
d'environ 300 images acquises pendant le mouvement du bras, jouant le rôle d'une vidéo, mais stockée
physiquement comme des fichiers DICOM individuels (un fichier par image, sans tag DICOM natif de
vidéo/temporalité).

\section{Objectif de la thèse}

L'objectif \textbf{n'est pas} de maximiser la performance brute d'un modèle de segmentation, mais de
comprendre \emph{comment} un modèle de deep learning arrive à ses segmentations sur ce type de données.
L'\textbf{interprétabilité / explicabilité} est la contribution centrale de la thèse ; un modèle de
segmentation performant est un moyen, pas une fin en soi.

Six structures anatomiques sont ciblées :

\begin{itemize}[nosep]
    \item \textbf{CLAV} --- clavicule
    \item \textbf{MSC} --- muscle scalène
    \item \textbf{VSC} --- veine sous-clavière
    \item \textbf{ASC} --- artère sous-clavière
    \item \textbf{K1} --- première côte
    \item \textbf{PB} --- plexus brachial
\end{itemize}

Ces six structures ont été confirmées mutuellement exclusives (pas de recouvrement inter-pixel), ce
qui oriente le choix de sortie du modèle vers un softmax multi-classe plutôt qu'un sigmoïde
multi-label.

\section{Plan de travail en six phases}

\begin{description}[style=nextline]
    \item[Phase 0 --- Préparation des données] Finalisation du nettoyage identifié en exploration
        (rééchantillonnage, normalisation des noms de segments, gestion des masques 4D/multi-couches),
        construction d'un jeu de données (image, masque) aligné et consolidé, définition d'un split
        train/val/test au niveau patient (14 patients --- attention aux fuites de données).
    \item[Phase 1 --- Modèle baseline] U-Net 2D simple (volontairement non état de l'art), entraîné
        sur les 20 images annotées par série. Dice/Hausdorff par structure comme référence de sanité,
        pas comme objectif d'optimisation.
    \item[Phase 2 --- Boîte à outils XAI] Seg-Grad-CAM (saillance par structure), MC Dropout (cartes
        d'incertitude par structure), analyse par occlusion. Validation qualitative des outils
        eux-mêmes avant d'en exploiter les résultats.
    \item[Phase 3 --- Analyse spécifique au caractère dynamique (axe original)] Cohérence temporelle
        des explications (la saillance dérive-t-elle de façon lisse d'une image à l'autre, comme
        l'anatomie réelle, ou de façon erratique ?), courbe de compression costoclaviculaire
        (distance CLAV--K1 en fonction du temps) pour localiser les images cliniquement critiques,
        corrélation entre incertitude/saillance et distance à l'image annotée la plus proche.
    \item[Phase 4 --- Ancrage clinique] Curation d'environ 10 à 15 cas représentatifs, session de
        validation qualitative avec le radiologue superviseur comparant les explications du modèle
        au raisonnement clinique.
    \item[Phase 5 --- Extension optionnelle] Entraînement d'une seconde architecture (p.\,ex.
        attention U-Net) sur le même split pour comparer le comportement d'interprétabilité à
        performance égale ; éventuellement TCAV sur des concepts cliniques, influence functions /
        leave-one-out sur les 20 images annotées.
    \item[Phase 6 --- Synthèse] Table récapitulative par structure (Dice + incertitude + qualité des
        explications + avis clinicien), rédaction finale.
\end{description}

\section{État d'avancement}

\begin{itemize}[nosep]
    \item[\checkmark] Exploration des données (EDA) --- voir chapitre~\ref{chap:exploration}
    \item[$\square$] Phase 0 --- Préparation des données
    \item[$\square$] Phase 1 --- Modèle baseline
    \item[$\square$] Phase 2 --- Boîte à outils XAI
    \item[$\square$] Phase 3 --- Analyse dynamique
    \item[$\square$] Phase 4 --- Ancrage clinique
    \item[$\square$] Phase 5 --- Extension optionnelle
    \item[$\square$] Phase 6 --- Synthèse
\end{itemize}
